\clearpage\thispagestyle{empty}
\pdfbookmark[0]{Sources for the July 1807 Report on Vesta}{vesta1807july-sources}
\begin{center}\textbf{Sources for the July 1807 Report on Vesta}\end{center}
\noindent This AI translation follows Franz Xaver von Zach’s July 1807 report in \href{https://gdz.sub.uni-goettingen.de/id/PPN236007467}{Gauss’s \textit{Werke}, vol.~VI} (Göttingen, 1874), pp.~288--290. The report includes Carl Friedrich Gauss’s letter of 7 June 1807 and his comparison of the observations with the second orbital elements of Vesta. The first printing is article VII, \textit{Fortgesetzte Nachrichten über den neuen von Dr. Olbers in Bremen entdeckten Planeten Vesta}, in the \textit{Monatliche Correspondenz zur Beförderung der Erd- und Himmelskunde}, vol.~XVI, \href{https://www.digitale-sammlungen.de/en/view/bsb10538607?page=91}{pp.~83}--\href{https://www.digitale-sammlungen.de/en/view/bsb10538607?page=100}{92}. Gauss’s letter begins on \href{https://www.digitale-sammlungen.de/en/view/bsb10538607?page=92}{p.~84}. Its comparison table runs from \href{https://www.digitale-sammlungen.de/en/view/bsb10538607?page=93}{p.~85} through \href{https://www.digitale-sammlungen.de/en/view/bsb10538607?page=94}{p.~86} to \href{https://www.digitale-sammlungen.de/en/view/bsb10538607?page=95}{p.~87}, where his remarks on three declinations follow.
\par\medskip
The table in the collected edition has a column for local mean time. Its two headings for the differences retain the historical labels R.$-$B. and R$-$B. The collected edition ends the excerpt with a bracketed statement on meridian observations by Groombridge and Maskeline. The first printing gives the corresponding Greenwich observations and spells the latter name Maskelyne on \href{https://www.digitale-sammlungen.de/en/view/bsb10538607?page=97}{p.~89}. The journal scans are from the Bayerische Staatsbibliothek. Translation and typesetting produced by AI.
\clearpage
